
\section{Square integrability of the actual theta periods}
\label{sec:sr-theta-period-square-integrability}

We retain the coefficient-square lattice majorants proved in
Section~\ref{sec:theta-decay-proof}. This section draws their two-variable
consequence. The first and second quotient variables have different
roles here; square integrability of the final theta period does not
follow by squaring a separate pointwise bound in the first variable.

For decreasing nonnegative vectors \(s=(s_1,\ldots,s_q)\) and
\(t=(t_1,\ldots,t_y)\), set
\[
 \begin{split}
 \rho_n(s)&=\sum_{i=1}^q(n-2i+1)s_i,\\
 P_{n,d}(s,t)&=
    d\sum_i s_i+n\sum_jt_j
                    -2\sum_{i,j}\min(s_i,t_j),\\
 \Delta(s,t)&=\max_{0\le k\le q}
             \left(\sum_{i\le k}s_i-\sum_{j\le k}t_j\right).
 \end{split}
 \tag{Q1}\label{eq:sr-theta-square-chamber-data}
\]
In a sum involving \(t_j\) beyond its length we put \(t_j=0\).
Write \(D=[E:\mathbb Q]>0\).
The proved coefficient-square and lattice estimates, after absorbing
the compact factors of the two Siegel sets, give
\[
 |\theta(f;a_X(s)k,a_W(t)k')|
       \le C_f\exp\{DP_{n,d}(s,t)/2\}.
 \tag{Q2}\label{eq:sr-theta-square-all-bound}
\]
On the actual full-rank cutoff they give, for every \(M>0\), a possibly
different \(C_{f,M}\) such that
\[
 |\theta(f;a_X(s)k,a_W(t)k')|
   \le C_{f,M}
       \exp\{DP_{n,d}(s,t)/2-M\Delta(s,t)\}.
 \tag{Q3}\label{eq:sr-theta-square-full-bound}
\]
Indeed these are square roots of
\eqref{eq:theta-all-torus} and
\eqref{eq:theta-full-torus}, with the already proved
coefficient-square majorant. The order of the Schwartz bound is
chosen before its compact finite support. Arbitrarily large \(M\)
is therefore permitted for every fixed test vector.
Finite sums of pure vectors satisfy the same estimates by the
triangle inequality. The actual cutoff projection gives the same
finite-sum assertion on the small side.

The Haar majorants for the first and second Siegel sets are respectively
\[
       C e^{-D\rho_n(s)}\,ds,\qquad
       C'e^{-D\rho_d(t)}\,dt.
 \tag{Q4}\label{eq:sr-theta-square-haar}
\]
These are the same rational Witt frames and Haar densities used in
the convergence proof, so no comparison with a different height
or a different local Fourier matrix is involved.

\begin{srlemma}[A partial-sum inequality]
\label{lem:sr-theta-square-majorization}
Suppose \(2n\le d\), \(q,y\le Y\), and \(2Y+1<n\), where all these
parameters are integers. Put \(M_0=d+q+1\).
Then
\[
 \frac{P_{n,d}(s,t)}2-\rho_n(s)-M_0\Delta(s,t)
 \le \frac{\rho_d(t)}2-\frac14\sum_jt_j
                 -\frac1{4(q+1)}\sum_i s_i .
 \tag{Q5}\label{eq:sr-theta-square-exponent}
\]
\end{srlemma}

\begin{proof}
Put \(\delta=\Delta(s,t)\ge0\), and let
\(u=(t_1,\ldots,t_q)\), padded with zeros or truncated as necessary.
For \(0\le k\le q\),
\(\sum_{i\le k}s_i\le\sum_{i\le k}u_i+\delta\).
We first prove
\[
 \sum_{i,j\le q}\min(s_i,s_j)
      -\sum_{i,j\le q}\min(s_i,u_j)
       \le\sum_{j\le q}(q-j)u_j+q\delta .
 \tag{Q6}\label{eq:sr-theta-square-minimum-slack}
\]
The summand on the left is at most \((s_i-u_j)_+\), because
\(\min(s_i,s_j)\le s_i\).
For fixed \(j\), the set \(s_i>u_j\) is an initial interval of length
\(k\). Hence
\[
 \sum_i(s_i-u_j)_+
  =\sum_{i\le k}s_i-ku_j
  \le\sum_{i\le k}(u_i-u_j)+\delta
  \le\sum_{i<j}u_i+\delta.
\]
For the last inequality discard the nonpositive terms with \(i\ge j\)
and bound each term with \(i<j\) by \(u_i\).
Summing over \(j\) proves (Q6).
Also
\[
 \sum_{i,j\le q}\min(s_i,u_j)
      \le\sum_{i\le q,j\le y}\min(s_i,t_j),
 \qquad
 \sum_i s_i\le\sum_j u_j+\delta .
\]

Set \(\varepsilon=1/[4(q+1)]\) and \(c=d/2-n+\varepsilon\ge0\).
For decreasing \(s\),
\[
 \sum_{i,j}\min(s_i,s_j)=\sum_i(2i-1)s_i,
 \qquad
 \rho_n(s)=n\sum_i s_i-\sum_{i,j}\min(s_i,s_j).
\]
Substitute these identities and (Q6) into the left side of (Q5),
and add \(\varepsilon\sum s_i\).
The result is at most
\[
 \sum_{j\le q}(c+q-j)u_j
       +\frac n2\sum_jt_j
       +(c+q-M_0)\delta .
 \tag{Q7}\label{eq:sr-theta-square-intermediate}
\]
Here we used \(c\sum s_i\le c\sum u_j+c\delta\).
The last term in (Q7) is nonpositive, since \(c\le d+1\).
For \(j\le q\),
\[
 c+q-j\le
       \frac{d-2j+1}2-\frac n2-\frac14:
\]
the difference of right and left is
\(n/2-q+1/4-\varepsilon>0\), using \(n\ge2Y+2\) and \(q\le Y\).
The right coefficient is nonnegative also for \(q<j\le y\),
because \(d\ge2n\) and \(j\le Y<n/2\).
We may therefore bound (Q7) by
\[
 \sum_{j\le y}
       \left(\frac{d-2j+1}2-\frac n2-\frac14\right)t_j
          +\frac n2\sum_jt_j
     =\frac{\rho_d(t)}2-\frac14\sum_jt_j .
\]
Subtracting \(\varepsilon\sum_i s_i\) proves (Q5).
\end{proof}

\begin{srproposition}[Two ranges of square integrability]
\label{prop:sr-theta-period-L2}
Let \(n=\dim_E X\), \(d=\dim_E W\), and let \(q,y\) be their
rational Witt indices. The actual theta period on \(U(W)\) is
square integrable in either of the following ranges.

First, for every Schwartz vector in the large convergence range,
it suffices that
\[
          q=0\ \hbox{or}\ d+2q<2n,\qquad
          y=0\ \hbox{or}\ n+y<d .
 \tag{Q8}\label{eq:sr-theta-period-L2-large-range}
\]
Second, for every vector in the actual small full-rank cutoff,
it suffices that
\[
       2n\le d,\qquad
       q,y\le Y,\qquad 2Y+1<n
                     \quad\hbox{for some integer \(Y\)}.
 \tag{Q9}\label{eq:sr-theta-period-L2-small-range}
\]
These assertions concern the specified Haar theta periods; changing a
quotient Haar measure multiplies them by a fixed finite scalar.
\end{srproposition}

\begin{proof}
For the first assertion use
\(P_{n,d}(s,t)\le d\sum_i s_i+n\sum_jt_j\) in (Q2).
Integrate its right side in the first quotient using (Q4).
The resulting first chamber integral is
\[
 \int_{s_1\ge\cdots\ge s_q\ge0}
   \exp\left\{D\sum_{i=1}^q(d/2-n+2i-1)s_i\right\}ds.
 \tag{Q10}\label{eq:sr-theta-period-L2-first-integral}
\]
For a decreasing nonnegative chamber, an exponential
\(\exp(\sum c_i s_i)\) is integrable exactly when every nonempty
initial coefficient sum is negative: substitute
\(r_i=s_i-s_{i+1}\), \(s_{q+1}=0\).
In (Q10) that initial sum is
\(Dk(d/2-n+k)\), negative for \(1\le k\le q\) under (Q8).
For \(q=0\) the chamber has volume \(1\).
Thus on the second Siegel set
\[
        |I_X(f)(a_W(t)k')|
               \le C\exp\{Dn\sum_jt_j/2\}.
\]
Its squared integral, with (Q4), is bounded by a constant times the
exponential chamber integral with coefficients
\(D(n-d+2j-1)\). The initial sums are \(Dk(n-d+k)\),
negative under the second condition of (Q8). This proves the large
assertion without applying a divergent integral of \(|\theta|^2\)
in the first variable.

For the second assertion choose the minor gain in (Q3) with
\(M=DM_0\), and apply (Q5). On the two Siegel sets this gives
\[
 |\theta(f;a_X(s)k,a_W(t)k')|
 \le C
 \exp\left\{D\rho_n(s)-\frac{D}{4(q+1)}\sum_i s_i\right\}
 \exp\left\{\frac D2\rho_d(t)-\frac D4\sum_j t_j\right\}.
 \tag{Q11}\label{eq:sr-theta-period-L2-separated}
\]
The first Haar density cancels its \(\rho_n(s)\) term.
The remaining integral
\(\int\exp\{-D\sum_i s_i/[4(q+1)]\}\,ds\) is finite.
Therefore
\[
 |I_X(f)(a_W(t)k')|
       \le C'\exp\{D\rho_d(t)/2-D\sum_jt_j/4\}.
\]
After squaring, the second Haar density cancels the
\(D\rho_d(t)\) term. The remaining bound
\(\exp\{-D\sum_jt_j/2\}\) is integrable on its chamber.
All estimates are uniform on the compact factors of the two
Siegel sets. Reduction theory transfers the chamber bounds to
the actual quotient integrals, proving square integrability.
The measurability needed here follows either from the locally
uniform Schwartz sums or directly from the dominated first period.
\end{proof}

\begin{srcorollary}[The source square-integrability condition]
\label{cor:sr-theta-SqInt}
Every actual geometric triple in the two theta towers satisfies the
square-integrability condition for its large polynomial theta periods
and its small polynomial theta periods in the actual cutoff.
\end{srcorollary}

\begin{proof}
Put \(Y=p+q\), in the notation of the geometric triple.
Its exact inequalities are
\[
 2Y+1<(a-a')/2,\qquad 2Y+1<a'.
\]
The large first Witt index is at most \(p\le Y\), and the target
Witt index is at most \(Y\).
With \(n=a,d=a+a'\), the first inequality gives
\(d+2p<2a\), and the second gives \(a+Y<d\).
Thus (Q8) applies.
For the small side \(n=a'\), \(d=a+a'>2a'\);
both its first Witt index and the target Witt index are at most
\(Y\), and \(2Y+1<a'=n\). Thus (Q9) applies.
The inequalities in particular give \(a'\ge2\), so no unproved
rank-zero convergence endpoint is being used.
\end{proof}
